% !PRINTED 158
\weberpage{158}

\begin{center}
\S.\,38.

\gesperrt{Substitution der Verhältnisse.}
\end{center}

Wir haben schon oben gesehen, dass zwei mit einer dritten ähnliche Substitutionen gleicher Dimension unter einander ähnlich sind. Demnach können wir alle mit einander ähnlichen Substitutionen $n^{\mathrm{ter}}$ Dimension in eine Classe vereinigen, und jede Substitution kann in einer und nur in einer solchen Classe untergebracht werden. Die einzelnen Substitutionen einer Classe heissen die \gesperrt{Repräsentanten} der Classe, und jede Classe ist durch irgend einen ihrer Repräsentanten völlig bestimmt.

Ist $A$ ähnlich mit $A'$, $B$ ähnlich mit $B'$, so ist auch $AB$ ähnlich mit $A'B'$.

Bezeichnen wir also mit $\mathfrak{A}, \mathfrak{B}$ die Classen, in die $A, A'$ und $B, B'$ gehören, so gelangt man immer in dieselbe Classe, welchen Repräsentanten aus $\mathfrak{A}$ und aus $\mathfrak{B}$ man auch zusammensetzen mag. Diese Classe, die durch $AB$ oder $A'B'$ repräsentirt wird, nennen wir daher aus $\mathfrak{A}$ und $\mathfrak{B}$ zusammengesetzt und bezeichnen sie mit $\mathfrak{A}\mathfrak{B}$. Bei dieser Zusammensetzung gelten dieselben Regeln, wie bei der Zusammensetzung der Substitutionen selbst, und die Gesammtheit der Classen bildet also auch eine Gruppe.

Diese Auffassung der Substitutionen und ihrer Zusammensetzung ist immer dann zweckmässig, wenn es, wie z.\,B. in der projectiven Geometrie, nur auf die Verhältnisse der Variablen ankommt, und wir bezeichnen daher die Substitutionsclassen als \gesperrt{Substitutionen der Verhältnisse.}

Wir werden aber diese Substitutionen der Verhältnisse ebenso bezeichnen, wie die anderen Substitutionen, nämlich jede Classe durch einen Repräsentanten, wodurch nicht leicht ein Missverständniss entstehen wird. Den Repräsentanten kann man nach sehr verschiedenen Gesichtspunkten auswählen.

Oft empfiehlt es sich, ihn so anzunehmen, dass die Determinante $= 1$ ist. Dadurch ist aber bei einer Substitution $n^{\mathrm{ter}}$ Dimension die multiplicative Substitution $(\mu)$ nur bis auf eine $n^{\mathrm{te}}$ Einheitswurzel bestimmt. Sind die Repräsentanten so gewählt, so bleibt die Eigenschaft bei der Composition erhalten, und diese Repräsentanten der Classen bilden also auch unter sich eine Substitutionsgruppe.
